\scanpage{34}
Quelques rappels de résultats de Kostant:

a) Pour $\mu_h\subset(T)\subset G$ fixé, les droites cycliques \inserted{$D$} $1$-normalisées par $\mu_h$ (\inserted{i.e.,} avec car $1$) forment un espace homogène sous $T$, stabilisateur $\mu_h$, donc un espace homogène sous $\mu_h$.
Les droites cycliques $D$ correspondent biunivoquement aux sous tores $T'$ de $G$ \deleted{tels que}\inserted{en opposition} \deleted{$\unsure{T|\mu_h}$} à $T$ pour $\mu_h$ (i.e. normalisé par $\mu_h$, et tel que $\mu_h$ y admet un \deleted{\unsure{elt}}\inserted{car.} propre au moins qui est \deleted{\unsure{racine}} car. \uline{primitif}, ($\unsure{\simeq|\Leftrightarrow}$ le car. $1$ y est car propre): à $D$ correspond $T'=\mathrm{Cent}(D)$, à $T'$ correspond l'\deleted{unique droite de}\unsure{le sous-espace|le sous-espace propre|unique sous-espace} propre de $\mu_h$ opérant sur \deleted{$\unsure{\mathfrak t}$} $\mathrm{Lie}\,T'\subset\mathfrak g$. On a d'ailleurs $T\cap N(T')=\mu_h$.

b) Les droites cycliques $D$ de $\mathfrak g$ \uline{mod} $\mu_h^*$ s'identifient can. \unsure{avec|à|via} une droite $\Delta$ canonique sur $k$. On l'obtient par descente à partir de $D^{\otimes h}$ défini sur le schéma des droites cycliques qui devient un schéma en droites sur le schéma des tores maximaux $T'$, et à ce titre se descend.
% slanted
\margin{NB Si on se place sur un topos, tel le classifiant ${\unsure{\underline B|B}}_{E_k}^{\deleted{\unsure{*}}}$, on ne pourra pas en général trouver de $a$, donc on ne pourra définir de torseur $\mathrm{KostKillnorm}$: cela provient du fait que la repr. de $E$ sur $\Delta$ peut être non triviale. Il faut passer à \unsure{rev.|une} étale d'indice $2$ \unsure{---}. De façon générale, on a obstruction \deleted{\unsure{pour}}\inserted{à trouver $D$} dans $H^2(\unsure{Q|\mathbb Q};E;\mu_h)$ nulle si on restreint \unsure{à|:} $H^{\unsure{2|1}}(Q,\mu_h)$ ou $H^2(Q;E^\circ;\mu_h)$}
On l'obtient aussi comme l'image du lieu cyclique de $\mathfrak g$ par l'\unsure{hom|hom.} de Kostant $\mathfrak g\longrightarrow I=\mathfrak t/W$, qui a une structure \inserted{\uline{vectorielle}} \deleted{\unsure{linéaire}} car le stabilisateur \inserted{dans $W$} d'une droite cyclique $D\subset\mathfrak t$ y opère par $\mu_h$, et $D/\mu_h$ a une structure \deleted{de droite}\inserted{vectorielle} \ldots

Il serait intéressant de calculer \deleted{les relations de \unsure{cette}} la repr qu'on obtient ainsi de $\unsure{E|\mathcal E|G}$ (autom. ext.) dans $\mathbb Z/2\mathbb Z\simeq\mu_2(\mathbb Z)$, et la comparer avec le car. provenant de la repr de $E$ sur $\mathfrak g_\psi$ ($\psi$ plus grande racine). P. ex. il est clair par les résultats de Kostant que si \deleted{la} deuxième repr. est triviale, la première aussi.
A-t-on, pour tt tore $T$ dans un Borel \deleted{\unsure{g}}$B$, un isom can $\mathfrak g_\psi\simeq\Delta$ (ou $\mathfrak g_\psi\simeq\Delta^{-1}$, \deleted{\unsure{c}} cela revient au $=$, car \unsure{on a} un isom can $\Delta^{\otimes2}\simeq 1$).
% archive sheet number: 33
